% 사과 가격 예측 시뮬레이터 - 수학식 LaTeX 문서 (초등학생용)
% 문서 버전: v1.0.0.0
% 대상: 초등학생 (5-6학년)
% 작성일: 2026-02-02
% 컴파일: xelatex 파일명.tex

\documentclass[12pt,a4paper]{article}

% 한글 지원
\usepackage{fontspec}
\setmainfont{Noto Sans CJK KR}
\setsansfont{Noto Sans CJK KR}

% 컬러 이모지 (twemojis 패키지)
\usepackage{twemojis}

% 수학 패키지
\usepackage{amsmath}
\usepackage{amssymb}

% 색상
\usepackage{xcolor}
\definecolor{teal}{RGB}{0,131,143}
\definecolor{darkblue}{RGB}{21,101,192}
\definecolor{lightcyan}{RGB}{224,247,250}
\definecolor{lightblue}{RGB}{227,242,253}
\definecolor{lightyellow}{RGB}{255,249,196}
\definecolor{lightgreen}{RGB}{232,245,233}
\definecolor{lightpink}{RGB}{252,228,236}

% 박스
\usepackage{tcolorbox}
\tcbuselibrary{skins,breakable}

% 페이지 설정 (여백 넓게)
\usepackage[margin=2.5cm]{geometry}

% 줄간격
\usepackage{setspace}
\setstretch{1.3}

% 제목 스타일
\usepackage{titlesec}
\titleformat{\section}{\Large\bfseries\color{teal}}{\thesection}{1em}{}
\titleformat{\subsection}{\large\bfseries\color{darkblue}}{\thesubsection}{1em}{}

% 하이퍼링크
\usepackage{hyperref}
\hypersetup{colorlinks=true,linkcolor=teal,urlcolor=darkblue}

\title{\Huge\textbf{\texttwemoji{red apple} 사과 가격 맞추기!}\\[0.5em]
\Large 수학 공식 모음 \textbf{(초등학생용)}}
\author{문서 버전: v1.0.0.0}
\date{2026년 2월 2일}

\begin{document}

\maketitle

\begin{tcolorbox}[colback=lightpink,colframe=magenta!50!black,title=\textbf{\texttwemoji{waving hand} 이 문서는요...}]
\large
사과 가격 예측 시뮬레이터에서 사용하는 \textbf{수학 공식}을 쉽게 정리했어요!

어려운 공식도 있지만, \textbf{계산 예시}를 따라하면 이해할 수 있어요 \texttwemoji{flexed biceps}
\end{tcolorbox}

\tableofcontents
\newpage

%=============================================
\section{\texttwemoji{star} 품질 점수 계산하기}
%=============================================

사과의 ``좋은 정도''를 숫자로 나타내요!

\subsection{\texttwemoji{thinking face} 품질 점수란?}

\begin{tcolorbox}[colback=lightcyan,colframe=teal,title=\textbf{\texttwemoji{sparkles} 품질 점수 = 당도 점수 + 무게 점수}]
\Large
\begin{equation}
\boxed{\text{품질} = \text{당도점수} + \text{무게점수}}
\end{equation}
\normalsize

\begin{itemize}
    \item \textbf{당도}: 사과가 얼마나 달콤한지 (0~50점)
    \item \textbf{무게}: 사과가 얼마나 무거운지 (0~50점)
    \item \textbf{품질}: 둘을 합친 점수 (0~100점)
\end{itemize}
\end{tcolorbox}

\subsection{\texttwemoji{honey pot} 당도 점수 계산}

\begin{tcolorbox}[colback=lightblue,colframe=darkblue,title=\textbf{당도 점수 공식}]
\Large
\begin{equation}
\text{당도점수} = \frac{\text{당도} - 10}{8} \times 50
\end{equation}
\normalsize

\textbf{당도 범위}: 10 ~ 18 (Brix라는 단위를 써요)
\end{tcolorbox}

\subsubsection*{\texttwemoji{pencil} 함께 계산해봐요!}

\begin{tcolorbox}[colback=lightyellow,colframe=orange]
\textbf{문제}: 당도가 \textbf{14 Brix}인 사과의 당도점수는?

\large
\begin{align}
\text{당도점수} &= \frac{14 - 10}{8} \times 50 \\[0.3em]
&= \frac{4}{8} \times 50 \\[0.3em]
&= 0.5 \times 50 \\[0.3em]
&= \boxed{25\text{점}}
\end{align}
\normalsize
\end{tcolorbox}

\subsection{\texttwemoji{balance scale} 무게 점수 계산}

\begin{tcolorbox}[colback=lightblue,colframe=darkblue,title=\textbf{무게 점수 공식}]
\Large
\begin{equation}
\text{무게점수} = \frac{\text{무게} - 150}{200} \times 50
\end{equation}
\normalsize

\textbf{무게 범위}: 150g ~ 350g
\end{tcolorbox}

\subsubsection*{\texttwemoji{pencil} 함께 계산해봐요!}

\begin{tcolorbox}[colback=lightyellow,colframe=orange]
\textbf{문제}: 무게가 \textbf{250g}인 사과의 무게점수는?

\large
\begin{align}
\text{무게점수} &= \frac{250 - 150}{200} \times 50 \\[0.3em]
&= \frac{100}{200} \times 50 \\[0.3em]
&= 0.5 \times 50 \\[0.3em]
&= \boxed{25\text{점}}
\end{align}
\normalsize
\end{tcolorbox}

\subsection{\texttwemoji{party popper} 품질 점수 완성!}

\begin{tcolorbox}[colback=lightgreen,colframe=green!50!black,title=\textbf{종합 예시}]
\textbf{당도 14 Brix, 무게 250g인 사과}

\large
\begin{align}
\text{당도점수} &= 25\text{점} \\
\text{무게점수} &= 25\text{점} \\[0.3em]
\text{품질} &= 25 + 25 = \boxed{50\text{점}}
\end{align}
\normalsize

\texttwemoji{light bulb} 품질 50점은 ``중간 정도'' 좋은 사과예요!
\end{tcolorbox}

%=============================================
\section{\texttwemoji{crystal ball} 가격 예측하기}
%=============================================

품질 점수로 사과 가격을 맞춰봐요!

\subsection{\texttwemoji{brain} AI가 배우는 공식}

\begin{tcolorbox}[colback=lightcyan,colframe=teal,title=\textbf{\texttwemoji{robot} AI의 예측 공식}]
\Large
\begin{equation}
\boxed{\text{가격} = w \times \text{품질} + b}
\end{equation}
\normalsize

\begin{itemize}
    \item $w$ (가중치): 품질 1점이 가격에 얼마나 영향을 주는지
    \item $b$ (기본값): 품질이 0점일 때의 기본 가격
\end{itemize}
\end{tcolorbox}

\subsection{\texttwemoji{chart increasing} 학습 전 vs 학습 후}

\begin{center}
\large
\begin{tabular}{|c|c|c|c|}
\hline
& \textbf{$w$ (가중치)} & \textbf{$b$ (기본값)} & \textbf{예측 정확도} \\
\hline
\texttwemoji{hatching chick} 학습 전 & 아무 숫자 & 0 & \texttwemoji{cross mark} 엉터리 \\
\hline
\texttwemoji{eagle} 학습 후 & 약 30 & 약 1000 & \texttwemoji{check mark button} 정확! \\
\hline
\end{tabular}
\normalsize
\end{center}

\subsection{\texttwemoji{pencil} 가격 예측 계산}

\begin{tcolorbox}[colback=lightyellow,colframe=orange,title=\textbf{예시: 품질 50점 사과의 가격은?}]
학습된 AI: $w = 30$, $b = 1000$

\large
\begin{align}
\text{가격} &= 30 \times 50 + 1000 \\[0.3em]
&= 1500 + 1000 \\[0.3em]
&= \boxed{2500\text{원}}
\end{align}
\normalsize

\texttwemoji{light bulb} 품질 50점 사과는 약 2,500원이에요!
\end{tcolorbox}

%=============================================
\section{\texttwemoji{rocket} AI는 어떻게 배울까?}
%=============================================

\subsection{\texttwemoji{thinking face} 학습의 원리}

\begin{tcolorbox}[colback=lightpink,colframe=magenta!50!black]
\large
\begin{enumerate}
    \item AI가 처음에 \textbf{아무렇게나} 가격을 예측해요
    \item 실제 가격과 \textbf{얼마나 틀렸는지} 확인해요 (오차)
    \item 오차를 줄이는 방향으로 $w$와 $b$를 \textbf{조금씩 바꿔요}
    \item 이 과정을 \textbf{수백 번} 반복해요
    \item 점점 예측이 정확해져요! \texttwemoji{party popper}
\end{enumerate}
\normalsize
\end{tcolorbox}

\subsection{\texttwemoji{turtle} vs \texttwemoji{rabbit} 학습 속도}

\begin{center}
\large
\begin{tabular}{|c|c|c|}
\hline
\textbf{학습률} & \textbf{비유} & \textbf{결과} \\
\hline
너무 작으면 & \texttwemoji{turtle} 거북이 & 너무 느려요 \\
\hline
적당하면 & \texttwemoji{person walking} 사람 & \texttwemoji{check mark button} 딱 좋아요! \\
\hline
너무 크면 & \texttwemoji{rabbit} 토끼 & 목표를 지나쳐요 \texttwemoji{dizzy} \\
\hline
\end{tabular}
\normalsize
\end{center}

%=============================================
\section{\texttwemoji{trophy} 얼마나 잘 맞췄을까?}
%=============================================

\subsection{\texttwemoji{hundred points} R² 점수}

\begin{tcolorbox}[colback=lightcyan,colframe=teal,title=\textbf{R² 점수란?}]
\large
AI가 \textbf{얼마나 잘 예측했는지} 알려주는 점수예요!

\begin{itemize}
    \item \textbf{1에 가까우면}: 아주 잘 맞췄어요! \texttwemoji{star}
    \item \textbf{0에 가까우면}: 그냥 평균이랑 비슷해요
    \item \textbf{0보다 작으면}: 평균보다 못 맞췄어요 \texttwemoji{crying face}
\end{itemize}
\normalsize
\end{tcolorbox}

\subsection{\texttwemoji{magnifying glass tilted left} R² 점수 읽는 법}

\begin{center}
\large
\begin{tabular}{|c|c|c|}
\hline
\textbf{R² 점수} & \textbf{의미} & \textbf{상태} \\
\hline
0.8 이상 & 아주 잘 맞췄어요! & \texttwemoji{check mark button} 양호 \\
\hline
0.5 ~ 0.8 & 그럭저럭 맞췄어요 & \texttwemoji{warning} 보통 \\
\hline
0.5 이하 & 더 연습이 필요해요 & \texttwemoji{cross mark} 부족 \\
\hline
\end{tabular}
\normalsize
\end{center}

%=============================================
\section{\texttwemoji{light bulb} 핵심 정리}
%=============================================

\begin{tcolorbox}[colback=lightgreen,colframe=green!50!black,title=\textbf{\texttwemoji{memo} 오늘 배운 공식들}]
\large

\textbf{1. 품질 점수}
\begin{equation}
\text{품질} = \text{당도점수} + \text{무게점수}
\end{equation}

\textbf{2. 당도 점수}
\begin{equation}
\text{당도점수} = \frac{\text{당도} - 10}{8} \times 50
\end{equation}

\textbf{3. 무게 점수}
\begin{equation}
\text{무게점수} = \frac{\text{무게} - 150}{200} \times 50
\end{equation}

\textbf{4. 가격 예측}
\begin{equation}
\text{가격} = w \times \text{품질} + b
\end{equation}

\normalsize
\end{tcolorbox}

%=============================================
\section{\texttwemoji{books} 참고자료}
%=============================================

\begin{center}
\begin{tabular}{|c|l|l|}
\hline
\textbf{\#} & \textbf{자료명} & \textbf{설명} \\
\hline
1 & 사과 가격 예측 시뮬레이터 & 직접 실습해볼 수 있어요! \\
\hline
2 & 수업지도안 (초등학생용) & 수업 내용 정리 \\
\hline
\end{tabular}
\end{center}

%=============================================
\section{\texttwemoji{memo} 문서 정보}
%=============================================

\begin{center}
\begin{tabular}{|c|c|c|c|l|c|}
\hline
\textbf{버전} & \textbf{날짜} & \textbf{시간} & \textbf{변경} & \textbf{내용} & \textbf{작성자} \\
\hline
v1.0.0.0 & 2026-02-02 & 11:30 & 최초 & 초안 작성 (초등학생용) & Claude \\
\hline
\end{tabular}
\end{center}

\vspace{1cm}
\begin{center}
\Large
\textit{\texttwemoji{red apple} 사과 가격 맞추기 - 수학 공식 모음 (초등학생용)}

\vspace{0.5cm}
\normalsize
\texttwemoji{waving hand} 수고했어요! 이제 시뮬레이터로 직접 실험해봐요!
\end{center}

\end{document}
